% Source-derived chapter 14: Some classical irreducible rigid local systems
% Original source: rigid-local-systems-multiplicative-eigenvalue-problem
\chapter{Some classical irreducible rigid local systems}
\label{rigid-local-systems-multiplicative-eigenvalue-problem-chapter-14}
\source{rigid-local-systems-multiplicative-eigenvalue-problem}

\begin{remark}[Source section]
\label{rigid-local-systems-multiplicative-eigenvalue-problem-chapter-14-source}
\source{rigid-local-systems-multiplicative-eigenvalue-problem}
\source{rigid-local-systems-multiplicative-eigenvalue-problem}
This chapter reproduces the corresponding section of the retrieved
LaTeX source, including its definitions, statements, examples, and
proofs. It has not yet been translated into Lean.
\notready
\end{remark}

% The source section title was: Some classical irreducible rigid local systems
\label{classH}
\source{rigid-local-systems-multiplicative-eigenvalue-problem}
The hypergeometric functions ${}_{\ell}F_{\ell-1}$ which give irreducible rigid rank $\ell$ local systems on $\Bbb{P}^1-\{0,1,\infty\}$, and the Pochhammer hypergeometric functions, which give irreducible rank $\ell$ local systems on $\Bbb{P}^1-\{\ell+1\text{ points}\}$, are classically well-studied rigid local systems. We recall \cite{BH, Haraoka} criteria for when these are unitary. The corresponding strange duals (Theorem \ref{egregium}) will give us F-vertices of suitable $P_k(3)$ and $P_k(\ell+1)$ respectively. This produces a rich collection of vertices in the multiplicative eigenvalue problem.
%Note that $n$ is now the rank of the rigid irreducible local system, and not $\ell$.
\subsection{The hypergeometric functions ${}_{\ell}F_{\ell-1}$}\label{hyper}
\source{rigid-local-systems-multiplicative-eigenvalue-problem}
Let $\alpha_1,\dots,\alpha_{\ell};\beta_1,\dots,\beta_{\ell}\in \Bbb{C}$. Let $D=z\frac{d}{dz}$ and consider the differential equation
\begin{equation}\label{math383}
\source{rigid-local-systems-multiplicative-eigenvalue-problem}
z(D+\alpha_1)\cdots (D+\alpha_{\ell})F=(D+\beta_1-1)\cdots (D+\beta_{\ell}-1)F
\end{equation}
By classical theory this is a differential equation with regular singularities, which are at $0,1,\infty$.
The local exponents are $1-\beta_1,\dots,1-\beta_{\ell}$ at $z=0$, $\alpha_1,\dots,\alpha_{\ell}$, at $z=\infty$ and
$0,1,\dots,\ell-2, -1+ \sum_{i=1}^{\ell}(\beta_i-\alpha_i)$ at $z=1$ (the local monodromies are $\exp(2\pi\sqrt{-1}x)$ where $x$ runs through the local exponents).
\begin{defi}
\source{rigid-local-systems-multiplicative-eigenvalue-problem}
\notready
For $\alpha\in \Bbb{R}$ define its fractional part $\langle \alpha\rangle$ by $\langle \alpha\rangle=\alpha-\left \lfloor{\alpha}\right \rfloor $.
\end{defi}

\begin{theorem}
\source{rigid-local-systems-multiplicative-eigenvalue-problem}
\notready\label{BHhere}\cite[Corollary 4.7]{BH}
\source{rigid-local-systems-multiplicative-eigenvalue-problem}
The monodromy of the differential equation \eqref{math383} can be conjugated to the unitary group $U(\ell)$ if and only if the $\alpha_i$ and $\beta_i$ are real, and
either of the following interlacing conditions hold (after rearranging the $\alpha_i$ and $\beta_i$ so that their fractional parts are weakly increasing):
$$\langle\alpha_1\rangle <\langle\beta_1\rangle<\langle\alpha_2\rangle <\langle\beta_2\rangle<\dots<\langle\alpha_{\ell}\rangle <\langle\beta_{\ell}\rangle,\ \ \ \text {or }\ \langle\beta_1\rangle<\langle\alpha_1\rangle <\langle\beta_2\rangle<\langle\alpha_2\rangle <\dots<\langle\beta_{\ell}\rangle<\langle\alpha_{\ell}\rangle .$$
\end{theorem}
\iffalse
\subsubsection{An example}\label{nmk}
\source{rigid-local-systems-multiplicative-eigenvalue-problem}
Take $\ell=2$ and $(\alpha_1,\alpha_2)=(\frac{1}{4},\frac{3}{4})$ and $(\beta_1,\beta_2)=(\frac{1}{2},1)$.  The corresponding local system has exponents $\frac{1}{2},0$ at $z=0$, $\frac{1}{4},\frac{3}{4}$, at $z=\infty$ and
$0, -\frac{1}{2}$ at $z=1$.  Using this data as $\mathcal{A}$ in Theorem \ref{egregium}, we compute $v(\ma)$, an F-vertex of $P_4(3)$: It has components $(\frac{1}{2},0,0,-\frac{1}{2})$ and $(\frac{1}{4},\frac{1}{4},-\frac{1}{4},-\frac{1}{4})$. Multiplying the second and third conjugacy classes by $i$ and $-i$ respectively, we find  same  vertex as the one in Section \ref{oldie}.
\fi
\subsection{The Pochhammer equation}\label{pokhie}
\source{rigid-local-systems-multiplicative-eigenvalue-problem}
Let $t_1,\dots,t_{\ell}$ be $\ell>1$ distinct points in $\Bbb{A}^1=\Bbb{P}^1-\{\infty\}$, and let $\lambda_1,\dots,\lambda_{\ell},\rho$ be complex numbers
satisfying $\sum_{i=1}^{\ell}\lambda_i\neq \ell\rho$. Let

$$A(\lambda,\rho)=
  \left[ {\begin{array}{cccc}
   \lambda_1 & \lambda_1-\rho & \dots & \lambda_1-\rho \\
   \lambda_2-\rho & \lambda_2 &\dots  & \lambda_2-\rho\\
   \vdots & \vdots &\vdots & \vdots \\
   \lambda_{\ell}-\rho & \lambda_{\ell}-\rho & \dots & \lambda_{\ell}\\
\end{array} }\right],\  T= \left[ {\begin{array}{cccc}
   t_1 & 0 & \dots & 0 \\
   0 & t_2 &\dots  & 0\\
   \vdots & \vdots &\vdots & \vdots \\
   0 & 0 & \dots & t_{\ell}\\\end{array} }\right].$$

The system of equations, with $Y=(y_1(z),\dots,y_{\ell}(z))^T$
\begin{equation}\label{math383sys}
\source{rigid-local-systems-multiplicative-eigenvalue-problem}
(z-T)\frac{dY}{dz}= A(\lambda,\rho)Y
\end{equation}
is the Pochhammer system of rank $\ell$, and is known to have regular singularities. The exponents at $t_i$ are $0$ with multiplicity $\ell-1$ and $\lambda_i$ with multiplicity $1$, and
the local exponents at $\infty$ are $-\rho$ with multiplicity $\ell-1$ and $-\rho'$ of multiplicity $1$ where
$\rho'=\sum_{i=1}^{\ell}\lambda_i -(\ell-1)\rho.$



The corresponding local system is irreducible if and only if
$\lambda_1-\rho,\lambda_2-\rho,\dots,\lambda_{\ell}-\rho,\rho,\rho'$ are all not in $\Bbb{Z}$ \cite[Proposition 1.3]{Haraoka}. We will assume that this is indeed the case. In this case the monodromy group can be conjugated to $U(\ell)$, by a result of Haraoka, if and only if one of the following conditions holds \cite[Proposition 1.5]{Haraoka}
\begin{enumerate}
\item $\langle\rho\rangle <\langle\lambda_j\rangle, j=1,\dots,\ell$ and $\sum_{j=1}^{\ell} \langle\lambda_j\rangle<(\ell-1)\langle\rho\rangle +1$, or
\item $\langle\rho\rangle >\langle\lambda_j\rangle, j=1,\dots,\ell$ and $(\ell-1) \langle\rho\rangle< \sum_{j=1}^{\ell} \langle\lambda_j\rangle$.
\end{enumerate}
\iffalse
\subsubsection{An example}
Let $\ell=3$, $\lambda_1=\lambda_2=\lambda_3=\frac{1}{3}$ and $\rho= \frac{1}{6}$. The first of the conditions
above is satisfied, and hence the monodromy is unitary. Now $\rho'=1-2\rho=\frac{2}{3}$. The local monodromies in Theorem \ref{egregium} (i.e., the exponents) are $(0,0,\frac{2}{6})$ at $t_1$, $t_2$ and $t_3$ and $(-\frac{1}{6},-\frac{1}{6},-\frac{4}{6})$ at $\infty$. The conjugacy class data is $n=6$, $\mu^1=\mu^2=\mu^3=(2,0,0)$
and $\mu^4=(5,5,2)$. This gives $\lambda^1=\lambda^2=\lambda^3=(1,1,0,0,0,0)=\omega_2$, and $\lambda^4 = (3,3,2,2,2,0)=\omega_2+2\omega_5$ at level $3$. Now $\kappa(\omega_2/3)=(1/3,1/3,0,0,0,0)-1/9(1,1,1,1,1,1)= (2/9,2/9,-1/9,-1/9,-1/9,-1/9)$ and $\kappa(\lambda^4/3)= (1,1,2/3,2/3,2/3,0)-2/3(1,1,1,1,1,1)=(1/3,1/3,0,0,0,-2/3)$.

 The point $(\kappa(\lambda^1/3),\kappa(\lambda^2/3),\kappa(\lambda^3/3),\kappa(\lambda^4/3))$ is an F-vertex of $P_6(4)$.
 \fi
 \iffalse
\begin{remark}
\source{rigid-local-systems-multiplicative-eigenvalue-problem}
\notready\label{pokeman}
\source{rigid-local-systems-multiplicative-eigenvalue-problem}
Consider a Pochhammer rigid local system  whose local monodromies are $n$th roots of unity.
Write $\rho=\tilde{\rho}/n+\left \lfloor{\rho}\right \rfloor $ and $\lambda_j=\tilde{\lambda}_j/n +\left \lfloor{\lambda_j}\right \rfloor$, with integers  $0<\tilde{\rho}<n$ (since $\rho\not\in\Bbb{Z}$) and $0\leq \tilde{\lambda}_j<n$.

Assume we are in the first case above. We then have $\tilde{\lambda}_j=\tilde{\rho}+u_j$ for integers $u_j>0$, and  the second condition  becomes $\sum_{j=1}^{\ell} u_j<n-\tilde{\rho}$. Since $\ell\leq\sum_{j=1}^{\ell} u_j$, we get that
$\ell\leq n-2$. If we are in the second case, we reach the same conclusion: Write $\tilde{\lambda}_j=\tilde{\rho}- u_j$ for integers $u_j>0$. The second condition becomes $0< \tilde{\rho}-\sum u_j$, or that $\sum_{j=1}^{\ell} u_j<\tilde{\rho}\leq n-1$, so that $\ell\leq n-2$ again.

The rank $\ell$ of a Pochhammer unitary rigid local system is therefore no more than $n-2$, assuming that all local monodromies are $n$th roots of unity. The condition that the  rigid local system is unitary is essential for this conclusion. For the hypergeometric local systems, using \cite[Corollary 4.7] {BH}  we reach a similar conclusion: The rank of a hypergeometric  unitary rigid local system is no more than $n/2$, assuming that all local monodromies are $n$th roots of unity (see Remark \ref{zoom}).
\end{remark}
\fi
